\documentclass{situationist-preprint}

\title{The Back of the Cube}
\msid{PHE-03-b}
\msfamily{PHE-03}
\msdiscipline{Phenomenology}
\mstype{Companion manuscript}
\mskeywords{adumbration; cube; horizon; empty intention; sensorimotor theory}

\begin{document}
\maketitle

\begin{abstract}
The cube is the standard example of perception in adumbrations: seen from the front, its back is not seen but co-intended. The example invites an objection to PHE-03-a. A cube has six faces, a finite number, so a short exploration should bring all of them to givenness, and convergence on complete givenness should be possible at least here. We give a detailed analysis of the back of the cube at three levels: faces, surfaces and interior. At the level of faces, at most three are seen from any point, but all six can be brought to fullness in turn. At the level of surfaces and interior, the back is not a fixed region but reappears at every grain and behind every cut. The cube confirms the result of PHE-03-a.
\end{abstract}

\section{The example and the objection}

Husserl's analyses of the spatial thing in \emph{Thing and Space} and \emph{Ideas I} return repeatedly to the cube or die \citep{Husserl1973,Husserl1913}. Seen from one standpoint, the cube shows a front face, perhaps two side faces, and its edges in perspectival foreshortening. The back is not seen, but it is not absent from the perception either: the cube is seen as a cube, with a back, and the back is intended emptily, with a determinable indeterminacy that further perception would fill. This is the structure of adumbration that PHE-03-a states in Definition~3.1 and Principles~P1--P2.

The example invites an objection to PHE-03-a. Its main result (Theorem~4.4) depends on Principle~P3, that a spatial thing has infinitely many determinations, together with Principle~P4, that a single phase of perception can hold only boundedly many determinations in fullness. But a cube is a paradigm of finite structure. It has six faces, twelve edges and eight vertices. If complete givenness of the cube is givenness of these, a short exploration that turns the cube within the retentional horizon would seem to give all of it in fullness. We analyze the back at three levels to show where the objection goes wrong.

\section{Faces}

\begin{proposition}\label{p:three}
From no point outside an opaque cube are more than three of its faces seen.
\end{proposition}

\begin{proof}
A face of a convex opaque solid is seen from a point $p$ only if $p$ lies in the open half-space bounded by the plane of the face on the side away from the solid. For each of the three pairs of opposite faces of a cube, the two half-spaces are disjoint, since the planes are parallel and the solid lies between them. Hence $p$ lies in at most one half-space of each pair, and at most three faces are seen.
\end{proof}

Proposition~\ref{p:three} is the elementary geometric core of the claim that the back is never seen together with the front. It holds for every standpoint, so no position of the perceiver removes the back; a move only changes which faces are back faces. It also makes precise the first fact that PHE-03-a combines in Principle~P4: what is presented from one orientation excludes what is presented from the opposed orientation.

\begin{proposition}\label{p:six}
Each face of the cube can be given in fullness, and all six faces can be given in fullness within one course of exploration.
\end{proposition}

\begin{proof}
By Principle~P2 of PHE-03-a, each face is reachable by a course of movement; turning the cube twice through a right angle about two axes brings every face to the front. This is Lemma~4.2 of PHE-03-a applied to six determinations.
\end{proof}

Here the objection seems to win. The six faces can all be given in fullness, and if the turning is quick, the earlier faces are still retained when the last one appears. But two points must be noted. First, a face that has turned away is retained as just seen, not seen; in retention it is held as having been given, and its fullness fades within the retentional horizon (PHE-05). Second, and decisively, \enquote{all six faces have been seen} is not \enquote{the cube has been completely given}. The faces of the cube are six determinations among many, and the objection's finitude is the finitude of a geometric description, not of the thing.

\section{Surfaces and interior}

The back of a face is not only the face opposite it. Consider the front face itself, seen in fullness. At ordinary distance its color appears uniform and its surface smooth. Moving closer gives the grain of the wood or the texture of the paint, and what was given in fullness at the earlier distance turns out to have been given as smooth only at that scale. The finer grain was co-intended as a determinable indeterminacy: the surface was seen as having some texture or other. At each scale there is a \enquote{back} in the relevant sense, a determination that is co-intended and not given, and the course of movement that would give it is specified, as Husserl's kinaesthetic analyses require \citep{Husserl1973}. The same holds for color under illumination: each change of light gives the face under a further aspect, and no light is the one under which the color is given once and for all.

\begin{proposition}\label{p:interior}
The interior of a solid cube is never given as interior.
\end{proposition}

\begin{proof}
Whatever is seen of a solid is seen as surface, by the half-space condition used in Proposition~\ref{p:three}. To see the interior one must cut the cube, and the cut produces new faces, which are surfaces of the parts, with their own backs. The interior is therefore given only by transforming it into exterior, and as interior it remains co-intended.
\end{proof}

Proposition~\ref{p:interior} shows that the back of the cube is not a region that exploration exhausts. It is a structural feature of every presentation of the thing: a determination given in fullness at one phase is given with a horizon of further determinations, among them the back side, the finer grain, the light not yet tried, and the inside. This is Principle~P3 of PHE-03-a seen in the example that seemed to refute it.

\section{Objections}

\begin{objection}
The sensorimotor theory holds that the back of the cube is perceptually present, as accessible through movement \citep{Noe2004}. If the back is present, the front and back are given together after all.
\end{objection}

\begin{reply}
Presence as access is what Husserl calls co-intention, and the sensorimotor theory supplies a good account of it. But access is not fullness. The back is present as something that would appear if one moved; it is not present as it would then appear. PHE-03-a's result concerns fullness, and the distinction between being present as accessible and being given in fullness is the distinction between the empty and the fulfilled intention.
\end{reply}

\begin{objection}
A glass cube shows all six faces at once.
\end{objection}

\begin{reply}
It shows the back faces through the front, inverted in orientation and altered by refraction. They are given, but under a different sense from the one under which they would be given if faced, and the glass has determinations, such as its index of refraction and its internal flaws, that are open in exactly the way described above. Transparency changes what is co-intended; it does not abolish co-intention.
\end{reply}

\section{Conclusion}

The cube is finite as a geometric figure and open as a thing. At the level of faces, the back is never seen with the front, though each face can be seen in turn. At the level of surfaces and interior, the back reappears at every grain and behind every cut. We can return to the cube itself, and every side of it can be given; but the cube is never given all at once, and the series of its adumbrations does not converge, as Theorem~4.4 of PHE-03-a states.

\begin{thebibliography}{99}\small
\bibitem[Husserl(1913)]{Husserl1913} Husserl, E. (1913). Ideen zu einer reinen Phänomenologie und phänomenologischen Philosophie. Erstes Buch. \emph{Jahrbuch für Philosophie und phänomenologische Forschung} 1.
\bibitem[Husserl(1973)]{Husserl1973} Husserl, E. (1973). \emph{Ding und Raum. Vorlesungen 1907}. Ed. U.~Claesges. Husserliana XVI. Martinus Nijhoff.
\bibitem[No\"e(2004)]{Noe2004} No\"e, A. (2004). \emph{Action in Perception}. MIT Press.
\end{thebibliography}

\end{document}
